\documentclass[11pt,oneside]{article}
\usepackage[a4paper,margin=27mm,headheight=15pt]{geometry}
\usepackage[T1]{fontenc}
\usepackage[utf8]{inputenc}
\IfFileExists{libertinus.sty}{\usepackage{libertinus}}{\usepackage{lmodern}}
\usepackage{microtype}
\usepackage{amsmath,amssymb,amsthm,mathtools,mathrsfs}
\usepackage{bm}
\usepackage{booktabs,longtable,array,tabularx,multirow}
\newcolumntype{Y}{>{\raggedright\arraybackslash}X}
\usepackage{graphicx}
\usepackage{xcolor}
\usepackage{enumitem}
\usepackage{fancyhdr}
\usepackage{titlesec}
\usepackage{listings}
\usepackage{xurl}
\usepackage{hyperref}
\usepackage[nameinlink,noabbrev]{cleveref}

\definecolor{linkblue}{RGB}{25,75,135}
\definecolor{shadegray}{RGB}{246,247,249}
\definecolor{rulegray}{RGB}{195,201,208}
\definecolor{newgreen}{RGB}{20,105,66}
\definecolor{importblue}{RGB}{36,79,128}
\definecolor{conjectureorange}{RGB}{151,83,15}
\hypersetup{
 colorlinks=true,linkcolor=linkblue,citecolor=linkblue,urlcolor=linkblue,
 pdftitle={Half-Integer Spectral Arithmetic for the Fabius--Rvachev System},
 pdfsubject={Weighted Prouhet cancellation, dyadic q-binomial aliasing, cyclotomic traces, and a Lambert-W bridge},
 pdfkeywords={Fabius function,Rvachev up-function,Thue-Morse,q-binomial,sinc product,Lambert W,Prouhet,cyclotomic field,discrete cosine transform}}

\pagestyle{fancy}
\fancyhf{}
\fancyhead[L]{\small Half-integer spectral arithmetic}
\fancyhead[R]{\small\nouppercase{\leftmark}}
\fancyfoot[C]{\thepage}
\renewcommand{\headrulewidth}{0.4pt}
\titleformat{\section}{\Large\bfseries}{\thesection}{0.7em}{}
\titleformat{\subsection}{\large\bfseries}{\thesubsection}{0.7em}{}
\titleformat{\subsubsection}{\normalsize\bfseries}{\thesubsubsection}{0.7em}{}
\setcounter{secnumdepth}{3}
\setcounter{tocdepth}{2}
\setlist[itemize]{leftmargin=2em,itemsep=0.25em,topsep=0.4em}
\setlist[enumerate]{leftmargin=2.3em,itemsep=0.25em,topsep=0.4em}
\setlength{\emergencystretch}{3em}

\newtheorem{theorem}{Theorem}[section]
\newtheorem{proposition}[theorem]{Proposition}
\newtheorem{lemma}[theorem]{Lemma}
\newtheorem{corollary}[theorem]{Corollary}
\newtheorem{conjecture}[theorem]{Conjecture}
\newtheorem{problem}[theorem]{Problem}
\theoremstyle{definition}
\newtheorem{definition}[theorem]{Definition}
\newtheorem{algorithm}[theorem]{Algorithm}
\newtheorem{example}[theorem]{Example}
\theoremstyle{remark}
\newtheorem{remark}[theorem]{Remark}
\newtheorem{warning}[theorem]{Warning}

\newcommand{\R}{\mathbb R}
\newcommand{\C}{\mathbb C}
\newcommand{\Q}{\mathbb Q}
\newcommand{\N}{\mathbb N}
\newcommand{\Z}{\mathbb Z}
\newcommand{\Up}{\operatorname{up}}
\newcommand{\sinc}{\operatorname{sinc}}
\newcommand{\ord}{\operatorname{ord}}
\newcommand{\Gal}{\operatorname{Gal}}
\newcommand{\Tr}{\operatorname{Tr}}
\newcommand{\Norm}{\operatorname{N}}
\newcommand{\e}{\mathrm e}
\newcommand{\ii}{\mathrm i}
\newcommand{\dd}{\,\mathrm d}
\newcommand{\bigO}{\mathcal O}
\newcommand{\smallo}{o}
\newcommand{\eps}{\varepsilon}
\newcommand{\wt}{s_2}
\newcommand{\GF}{\mathcal F}
\newcommand{\PhiU}{\Phi}
\newcommand{\qbinom}[3]{\genfrac{[}{]}{0pt}{}{#1}{#2}_{#3}}
\newcommand{\newresult}{\textcolor{newgreen}{\textsc{New result proved here.}}}
\newcommand{\imported}{\textcolor{importblue}{\textsc{Imported baseline.}}}
\newcommand{\openstatus}{\textcolor{conjectureorange}{\textsc{Open or conjectural.}}}
\newenvironment{boxedremark}{\par\smallskip\noindent\begin{center}\begin{minipage}{0.94\textwidth}\hrule\vspace{0.6em}\small}{\vspace{0.6em}\hrule\end{minipage}\end{center}\smallskip}
\lstdefinestyle{pseudo}{basicstyle=\ttfamily\footnotesize,backgroundcolor=\color{shadegray},frame=single,
 rulecolor=\color{rulegray},columns=fullflexible,keepspaces=true,showstringspaces=false,breaklines=true,
 xleftmargin=0.7em,xrightmargin=0.7em,aboveskip=0.8em,belowskip=0.8em}

\title{\bfseries Half-Integer Spectral Arithmetic\\
for the Fabius--Rvachev System\\[0.35em]
\large Weighted Prouhet cancellation, dyadic $q$-binomial aliasing,\\
cyclotomic traces, and a Lambert--$W$ endpoint bridge}
\author{Frontier research report based on the recursive corpus\\
\small\path{Analysis/FabiusFunction/docs}\\
\small Vladimir Reshetnikov's \texttt{ProveIt} repository}
\date{27 August 2026}

\begin{document}
\maketitle

\begin{abstract}
Let
\[
 \PhiU(z):=\widehat{\Up}(z)
 =\prod_{j=0}^{\infty}\sinc\!\left(\frac z{2^j}\right),
 \qquad \sinc u:=\frac{\sin(\pi u)}{\pi u},
\]
be the Fourier transform of Rvachev's compactly supported up-function, and put
$a_k=\PhiU(k+\tfrac12)$.  The repository corpus already contains the Fourier
product, exact dyadic values of the Fabius function, their Gaussian-binomial and
Thue--Morse formulae, the Lambert--$W$ endpoint asymptotic, the complete real-axis
Fourier-decay theory assembled from several waves, and a recent spectral account
of the integer zero divisor.  This report does not relabel those results as new.
It develops instead the half-integer spectral sequence $(a_k)$, identified in the
previous spectral audit as a genuine unresolved bridge.

The first main theorem is an infinite, positive-weight Prouhet system of all
orders.  With $o_k=2k+1$, $\eps_k=(-1)^{s_2(k)}$, and
$w_k=|a_k|$, one has
\[
 \sum_{k\ge0}a_k=\frac12,
 \qquad
 \sum_{k\ge0}\eps_k w_k o_k^{2m}=0\quad(m\ge1).
\]
Thus the Thue--Morse partition of the odd integers has exactly matching positive
weighted even moments of every order.  More generally, the same coefficients
form an exact signed quadrature rule
$\sum_k\eps_kw_kP(o_k^2)=P(0)/2$ for every polynomial $P$.

The second main theorem is a finite dyadic spectral transform.  If $N=2^n$ and
$r$ is odd with $1\le r<N$, define
\[
 A_{n,r}=\frac1N\left[1-2\sum_{j=1}^{N-1}
 F\!\left(\frac jN\right)\cos\!\left(\frac{\pi rj}{N}\right)\right].
\]
Then the exact alias identity
\[
 A_{n,r}=\sum_{q\in\Z}\PhiU\!\left(qN+\frac r2\right)
\]
holds, the dyadic samples are recovered by an exact inverse discrete cosine
transform, and
\[
 |A_{n,r}-\PhiU(r/2)|\le
 \frac{2(1-2^{-n})\zeta(n)}{\pi^n}\,2^{-n(n-1)/2}
 \qquad(n\ge2).
\]
Substitution of the repository's finite $q=1/2$ Gaussian-binomial formula for
each $F(j/N)$ therefore gives superexponentially convergent, finite
$q$-binomial--Thue--Morse--cyclotomic approximants to every fixed half-integer
Fourier amplitude.

The folded coefficients form a Galois orbit in the maximal real
$2^{n+1}$-st cyclotomic field.  Their orbit polynomial lies in $\Q[X]$; all power
traces are explicit rational convolutions of dyadic Fabius values with a
Ramanujan-sum congruence filter.  These rational traces converge to nonlinear
spectral moments $\sum_k a_k^m$.  An additional exact factorization isolates a
finite Thue--Morse polynomial value as the sign carrier of every alias class.
Finally, the all-order Prouhet identities turn the endpoint function into an
exact ``beyond all algebraic orders'' spectral remainder.  Inserting the
repo's corrected Lambert--$W$ asymptotic produces a direct, explicit bridge
between the signed half-integer sinc amplitudes and the lower-Lambert saddle.
All new claims below are proved for human verification; none is represented as
already formalized in Lean.  Novelty is asserted only relative to the reviewed
repository corpus and targeted source searches, not as a universal priority
claim.
\end{abstract}

\begin{boxedremark}
\textbf{Status convention.}
\imported\ marks a theorem, identity, or notation already present in the
repository or in a cited source.  \newresult\ marks a statement for which a
proof is supplied here and which was not found in the audited TeX corpus or in
targeted searches of the cited literature.  \openstatus\ marks conjectures or
computational observations.  The distinction is essential because the report
is designed as a frontier supplement, not a rewrite of the existing corpus.
\end{boxedremark}

\tableofcontents
\newpage

\section{Scope, recursive audit, and the non-duplication boundary}\label{sec:scope}

\subsection{What was inspected}

The documentation tree is a stratified research record rather than one article.
The same-day recursive audit used for this report contains fifty-seven TeX
artifacts: twenty-five living or imported sources and thirty-two frozen historical
versions.  The living stratum includes the main Fabius--Rvachev exposition, a
mathematical glossary, the Lean walkthrough, the non-elementarity study, a
large consolidated non-formalized frontier volume, a Thue--Morse formula atlas,
a sequence of Fourier-decay reports and audits, and seven source-paper
transcriptions.  The historical stratum preserves earlier expositions, q-calculus
studies, asymptotic investigations, repeated-integration reports, gap registers,
and superseded formula atlases.  A path ledger is supplied in
\cref{app:inventory}.

The current canonical descendants were rechecked separately, because the
frontier directory is intentionally consolidated over time: processed drafts
may be absorbed into the large canonical volume or removed while their frozen
archive remains.  Accordingly, the appendix records the audited snapshot; it is
not a promise that every historical path remains a live branch-head file after
future repository maintenance.

Mathematical extraction was concentrated on the current descendants and on the
distinctive claims of each archived line of development.  Byte-identical and
superseded files were treated as provenance, not as independent evidence of
multiple results.  The audit was then compared with three frontier reports
produced from the same corpus on spectral determinants, spline/quadrature
structure, and arithmetic Fourier rays.  This prevents accidental duplication
of very recent deductions that are not yet part of the main exposition.

\subsection{Results explicitly treated as baseline}

The following are \emph{not} claimed as discoveries here.
\begin{enumerate}
 \item The relations between the bounded Fabius function $F$, its signed global
 extension $\GF$, and Rvachev's even compactly supported function $\Up$.
 \item The refinement differential equation, the entire Fourier product, the
 cosine-product form, and rapid decay of $\PhiU$.
 \item Exact rationality at dyadic points, the moment recurrences, the
 highest-bit recursion, and the finite $q=1/2$ Gaussian-binomial--Thue--Morse
 formula for $\GF(m/2^n)$.
 \item The integer zero set of $\PhiU$, multiplicity $1+\nu_2(m)$ at a nonzero
 integer $m$, strict logarithmic concavity between zeros, and the sign law
 $\operatorname{sgn}\PhiU(x)=(-1)^{s_2(M)}$ on $M<x<M+1$.
 \item The pointwise, metric, probabilistic, extremal, annular, and transfer-operator
 Fourier-decay results developed in the successive decay reports.
 \item The corrected endpoint asymptotic in the lower-Lambert coordinate,
 including the nonconstant Gamma--zeta periodic fluctuation.
 \item The spectral determinant model, spectral zeta and heat trace, all-order
 integer-zero jets, dyadic q-Euler determinants, and symmetric q-deformation
 developed in the prior spectral frontier report.
 \item Exact rational-ray decay laws, cyclotomic orbit norms, resultants,
 Ramanujan traces, and the Mersenne half-base bridge developed in the prior
 arithmetic-ray report.
\end{enumerate}

\subsection{The gap selected for this report}

The previous spectral report ends with the focused problem of classifying the
half-integer amplitudes
\[
 \PhiU\!\left(\frac{2k+1}{2}\right),
\]
which govern the leading Taylor coefficients at every integer zero.  The main
repository exposition contains the exact cosine series in which the same
amplitudes reconstruct $\Up$, while its dyadic arithmetic chapters contain
finite $q$-binomial formulae for the values being reconstructed.  Yet the
reviewed corpus does not develop the finite transform between those two arrays,
does not state the all-order weighted Prouhet identities below, and does not
connect those cancellations to the corrected Lambert endpoint scale.

This is therefore a natural unfilled intersection of five mature themes already
present separately in the repository:
\[
\begin{aligned}
 \text{Thue--Morse signs}
 &\longleftrightarrow \text{half-integer sinc amplitudes}\\
 &\longleftrightarrow \text{dyadic }q\text{-binomial values}
 \longleftrightarrow \text{cyclotomic traces}
 \longleftrightarrow \text{Lambert--}W.
\end{aligned}
\]

\subsection{Contribution map}

\cref{tab:contribution-map} summarizes the boundary between imported material,
new deductions proved here, and conjectural extensions.

{\small
\begin{longtable}{@{}>{\raggedright\arraybackslash}p{0.22\textwidth}>{\raggedright\arraybackslash}p{0.51\textwidth}>{\raggedright\arraybackslash}p{0.18\textwidth}@{}}
\caption{Boundary between imported material and the new theorem package.}\label{tab:contribution-map}\\
\toprule
Result & Content & Status \\
\midrule
\endfirsthead
\multicolumn{3}{c}{\small\itshape Contribution map (continued)}\\[0.3em]
\toprule
Result & Content & Status \\
\midrule
\endhead
\midrule
\multicolumn{3}{r}{\small\itshape Continued on the next page}\\
\endfoot
\bottomrule
\endlastfoot
Half-integer cosine spectrum & Isolates the sequence $a_k=\PhiU(k+1/2)$ as the nonzero Fourier coefficients of the two-periodic up-function. & \imported\ identity, new organization \\
Weighted Prouhet theorem & Proves equality of every positive weighted even moment in the Thue--Morse partition of the odd integers. & \newresult \\
Spectral quadrature & Shows $\sum_k\eps_kw_kP((2k+1)^2)=P(0)/2$ for every polynomial $P$. & \newresult \\
Dyadic moment filtration & Converts every derivative jet at a reduced dyadic point into a half-integer spectral moment and identifies its exact terminating depth. & \newresult\ formulation \\
Beyond-all-orders remainder & Rewrites $F(x)$ after subtraction of an arbitrary cosine Taylor polynomial, with all algebraic terms cancelling exactly. & \newresult \\
Exact dyadic alias transform & Gives a finite DCT of dyadic Fabius values equal to a complete folded half-integer alias class. & \newresult \\
Certified convergence & Gives a uniform $2^{-n^2/2+\bigO(n)}$ error bound from the folded coefficient to $\PhiU(r/2)$. & \newresult \\
$q$-binomial spectral approximants & Substitutes the finite Gaussian-binomial formula into the DCT, producing exact cyclotomic approximants. & \newresult \\
Cyclotomic orbit traces & Constructs rational orbit polynomials and explicit Ramanujan-filtered power traces. & \newresult \\
Spectral power-sum limits & Shows the rational Galois traces converge to $\sum_k a_k^m$ superexponentially up to a rapidly decreasing tail. & \newresult \\
Thue--Morse sign carrier & Factors each alias class into a finite root-of-unity Thue--Morse product and an absolutely convergent shifted spectral sum. & \newresult \\
Lambert bridge & Inserts the corrected lower-Lambert asymptotic into the exact Prouhet remainder. & new consequence of imported asymptotic \\
Alias-sign positivity & Conjectures positivity of the residual shifted spectral sum, equivalent to an exact Thue--Morse sign law for every folded coefficient. & \openstatus \\
\end{longtable}
}

\section{Baseline notation and identities}\label{sec:baseline}

\subsection{The three functions}

Let $F:[0,1]\to[0,1]$ be the bounded Fabius function.  Let
\begin{equation}\label{eq:global-F}
 \GF(x):=\sum_{n=0}^{\infty}\eps_n\Up(x-2n-1),
 \qquad \eps_n:=(-1)^{s_2(n)},
\end{equation}
be its signed global extension; the sum is locally finite.  On $[-1,1]$,
\begin{equation}\label{eq:Up-F-relations}
 \Up(t)=\GF(t+1)=
 \begin{cases}
  F(t+1),&-1\le t\le0,\\
  F(1-t)=1-F(t),&0\le t\le1.
 \end{cases}
\end{equation}
The function $\Up$ is even, nonnegative, $C^\infty$, supported on $[-1,1]$,
normalized by $\int_{\R}\Up=1$, and satisfies $\Up(0)=1$.
Every positive derivative of $\Up$ vanishes at $0$, because on the right of
zero it is $1-F(t)$ and $F$ is flat at the origin.

The global derivative identities are
\begin{equation}\label{eq:global-derivatives}
 \GF^{(m)}(x)=2^{\binom{m+1}{2}}\GF(2^m x),
 \qquad
 \Up^{(m)}(t)=2^{\binom{m+1}{2}}\GF\bigl(2^m(t+1)\bigr).
\end{equation}
At nonnegative integers,
\begin{equation}\label{eq:global-integers}
 \GF(2b)=0,
 \qquad
 \GF(2b+1)=\eps_b.
\end{equation}
These are repository baseline facts, ultimately descending from Arias de Reyna's
formulation \cite{ariascompact,ariasarith}.

\subsection{The Fourier product}

We use the Fourier convention
\[
 \widehat f(z)=\int_{\R}f(t)\e^{-2\pi\ii tz}\dd t.
\]
Put
\begin{equation}\label{eq:Phi-def}
 \PhiU(z):=\widehat\Up(z)
 =\prod_{j=0}^{\infty}\sinc\!\left(\frac z{2^j}\right),
 \qquad
 \sinc u=\frac{\sin(\pi u)}{\pi u}.
\end{equation}
The product is entire and even.  For every $J\ge1$ and real $x\ne0$,
selecting the first $J$ factors gives the elementary rapid-decay estimate
\begin{equation}\label{eq:Phi-power-bound-baseline}
 |\PhiU(x)|\le
 \frac{2^{J(J-1)/2}}{\pi^J|x|^J}.
\end{equation}
The nonzero real zeros are exactly the integers; at $m\ne0$ the multiplicity is
$1+\nu_2(|m|)$.  If $M<x<M+1$, then
\begin{equation}\label{eq:Phi-sign-law}
 \operatorname{sgn}\PhiU(x)=(-1)^{s_2(M)}.
\end{equation}

\subsection{The exact dyadic \texorpdfstring{$q$}{q}-binomial formula}

For later substitution, we record the current corpus formula.  With
$(q;q)_n=\prod_{j=1}^n(1-q^j)$ and the Gaussian binomial coefficient
$\qbinom{n}{k}{q}$, one has for $m,n\in\N$ and every rational translation $c$
\begin{equation}\label{eq:baseline-qbinomial}
\begin{aligned}
 \GF\!\left(\frac{m}{2^n}\right)
 ={}&\frac{1}{2^{n^2}(1/2;1/2)_n}
 \sum_{k=0}^{n}
 \frac{\qbinom{n}{k}{1/2}}
      {2^{k(k-1)}(n+k)!}\\
 &\qquad\times
 \sum_{s=0}^{m2^k-1}\eps_s
 \left(s-m2^k+c\right)^{n+k}.
\end{aligned}
\end{equation}
The right side is independent of $c$.  When $0\le m\le2^n$, it equals the
bounded value $F(m/2^n)$.  This is \imported\ from the main exposition; the
new use below is to insert it into a cyclotomic discrete cosine transform.

\subsection{The lower-Lambert endpoint theorem}

Write $L=\log2$ and, for sufficiently small $x>0$,
\begin{equation}\label{eq:lambda-W-def}
 W=W_{-1}(-Lx),
 \qquad
 \lambda(x)=-\frac WL,
 \qquad
 r(x)=2^{\lambda(x)}=\frac{\lambda(x)}x.
\end{equation}
Let $\Psi$ be the smooth one-periodic correction whose nonzero Fourier
coefficients are
\begin{equation}\label{eq:Psi-coeff-baseline}
 \widehat\Psi(k)=-\frac{\Gamma(-\chi_k)\zeta(1-\chi_k)}L,
 \qquad
 \chi_k=\frac{2\pi\ii k}{L}.
\end{equation}
Put
\begin{equation}\label{eq:g-constant}
 \mathfrak g=\gamma_1+\frac{\gamma^2}{2}-\frac{\pi^2}{12}.
\end{equation}
The corrected endpoint theorem in the corpus states
\begin{equation}\label{eq:Lambert-baseline}
 F(x)=
 \frac{\e^{\Psi(-W/L)}}{2^{7/12}\sqrt{\pi x}}
 \exp\!\left(\frac{\mathfrak g-W-\tfrac12W^2}{L}\right)
 \left(1+\bigO\!\left(\lambda(x)^{-1}\right)\right).
\end{equation}
The periodic factor is nonconstant and cannot be deleted from a sharp
asymptotic.  We shall not reprove \eqref{eq:Lambert-baseline}; instead we derive
an exact half-integer spectral identity whose left side is $F(x)$ and then
insert this imported theorem.

\section{The half-integer spectral sequence}\label{sec:half-spectrum}

\subsection{Definition and first values}

Define
\begin{equation}\label{eq:a-k-def}
 a_k:=\PhiU\!\left(k+\frac12\right),
 \qquad
 o_k:=2k+1,
 \qquad
 w_k:=|a_k|.
\end{equation}
The sign law \eqref{eq:Phi-sign-law} gives immediately
\begin{equation}\label{eq:a-sign}
 a_k=\eps_k w_k.
\end{equation}
The first amplitudes are shown in \cref{tab:first-amplitudes}.  Their rapid
shrinkage is deceptive: the exact identities below require cancellations over
all indices, and no finite initial segment has the moment-vanishing property.

\begin{table}[ht]
\centering\small
\begin{tabular}{@{}r r r@{}}
\toprule
$k$ & $\eps_k$ & $a_k=\PhiU(k+1/2)$ \\
\midrule
0 & $+1$ & $\phantom{-}0.553771275888810095344217416345$ \\
1 & $-1$ & $-0.046202299134306712992574862210$ \\
2 & $-1$ & $-0.008656867928610010895881610116$ \\
3 & $+1$ & $\phantom{-}0.001044808241827280615683472766$ \\
4 & $-1$ & $-0.000356152861629792122112229855$ \\
5 & $+1$ & $\phantom{-}0.000298368336838126419196617073$ \\
6 & $+1$ & $\phantom{-}0.000091404980598947955260616141$ \\
7 & $-1$ & $-0.000006555455730992891195331226$ \\
\bottomrule
\end{tabular}
\caption{Initial half-integer amplitudes.  The signs are the Thue--Morse word.}
\label{tab:first-amplitudes}
\end{table}

\subsection{Exact two-periodic cosine expansion}

Let $P$ denote the two-periodization of $\Up$,
\begin{equation}\label{eq:P-periodization}
 P(t):=\sum_{\ell\in\Z}\Up(t-2\ell).
\end{equation}
Because $\Up$ is supported on $[-1,1]$, $P(t)=\Up(t)$ for $-1\le t\le1$.

\begin{proposition}[Half-integer cosine spectrum]\label{prop:cosine-spectrum}
The Fourier series of $P$ is absolutely differentiable to every order and is
\begin{equation}\label{eq:cosine-spectrum}
 \boxed{
 P(t)=\frac12+\sum_{k=0}^{\infty}a_k\cos(\pi o_k t).}
\end{equation}
In particular, on $[-1,1]$ the series equals $\Up(t)$.
\end{proposition}

\begin{proof}
The $m$th complex Fourier coefficient for period two is
\[
 c_m=\frac12\int_{-1}^{1}\Up(t)\e^{-\pi\ii mt}\dd t
 =\frac12\PhiU(m/2).
\]
For nonzero even $m$, $m/2$ is a nonzero integer and $c_m=0$.  Also
$c_0=1/2$.  For $m=\pm(2k+1)$, evenness and reality give
$c_m=a_k/2$.  Thus complex Fourier inversion reduces to
\eqref{eq:cosine-spectrum}.  Estimate \eqref{eq:Phi-power-bound-baseline} with
arbitrarily large $J$ shows that $a_k$ is a Schwartz sequence, so the series
and every differentiated series converge absolutely and uniformly.
\end{proof}

\begin{remark}
The half-integer sequence is therefore not an auxiliary sampling choice.  It is
\emph{exactly} the nonconstant Fourier spectrum of the natural two-periodic
extension of $\Up$.  Integer Fourier samples disappear because the entire sinc
product vanishes there.
\end{remark}

\section{An all-order weighted Prouhet system}\label{sec:prouhet}

\subsection{Moment cancellation}

The classical Prouhet--Thue--Morse construction balances finitely many
\emph{unweighted} power sums on a finite dyadic block.  Here the half-integer
Fourier amplitudes supply positive, rapidly decreasing weights that balance
\emph{every} even power on the infinite set of odd integers.

\begin{theorem}[Half-integer weighted Prouhet theorem]\label{thm:weighted-prouhet}
The amplitudes \eqref{eq:a-k-def} satisfy
\begin{align}
 \sum_{k=0}^{\infty}a_k&=\frac12,
 \label{eq:mass-half}\\
 \sum_{k=0}^{\infty}a_k o_k^{2m}&=0
 \qquad(m\ge1).
 \label{eq:all-even-moments-zero}
\end{align}
Equivalently, for every $m\ge1$,
\begin{equation}\label{eq:positive-prouhet}
 \boxed{
 \sum_{\substack{k\ge0\\s_2(k)\ \mathrm{even}}}
 w_k(2k+1)^{2m}
 =
 \sum_{\substack{k\ge0\\s_2(k)\ \mathrm{odd}}}
 w_k(2k+1)^{2m}.}
\end{equation}
All displayed series converge absolutely.
\end{theorem}

\begin{proof}
Set $t=0$ in \eqref{eq:cosine-spectrum}.  Since $P(0)=\Up(0)=1$,
\[
 1=\frac12+\sum_{k\ge0}a_k,
\]
which proves \eqref{eq:mass-half}.  For $m\ge1$, termwise differentiation gives
\[
 P^{(2m)}(0)=(-1)^m\pi^{2m}
 \sum_{k\ge0}a_ko_k^{2m}.
\]
Every positive derivative of $P=\Up$ vanishes at zero, hence the left side is
zero.  This proves \eqref{eq:all-even-moments-zero}.  Absolute convergence
follows by applying \eqref{eq:Phi-power-bound-baseline} with $J>2m+1$.
Finally substitute $a_k=\eps_kw_k$ and move the negative-sign terms to the
opposite side.
\end{proof}

\begin{corollary}[Exact spectral quadrature at the origin]\label{cor:spectral-quadrature}
For every polynomial $Q\in\C[X]$,
\begin{equation}\label{eq:polynomial-quadrature}
 \boxed{
 \sum_{k=0}^{\infty}\eps_kw_k\,Q\bigl((2k+1)^2\bigr)
 =\frac{Q(0)}2.}
\end{equation}
Thus the signed positive weights $\eps_kw_k$ evaluate every polynomial in the
squared odd variable exactly at the off-grid point $0$.
\end{corollary}

\begin{proof}
Expand $Q(X)=\sum_{m=0}^d q_mX^m$ and use
\eqref{eq:mass-half}--\eqref{eq:all-even-moments-zero}.
\end{proof}

\begin{remark}[Why this is not the ordinary finite Prouhet theorem]
For a block $0\le k<2^n$, the unweighted Thue--Morse signs annihilate powers
only below degree $n$.  In \eqref{eq:positive-prouhet}, the support is infinite,
the nodes are the odd integers, the weights are the absolute half-integer sinc
amplitudes, and every even degree is annihilated.  The recent Thue--Morse
transform literature continues to develop finite Prouhet--Tarry--Escott
families \cite{cloitre2026}; it does not supply the spectral weights appearing
here.
\end{remark}

\subsection{A nontrivial signed measure with every positive moment zero}

\begin{corollary}[Spectral moment-null measure]\label{cor:moment-null-measure}
The finite signed measure
\begin{equation}\label{eq:nu-measure}
 \nu:=\frac12\sum_{k=0}^{\infty}a_k
 \bigl(\delta_{o_k}+\delta_{-o_k}\bigr)
\end{equation}
has total mass $1/2$ and satisfies
\begin{equation}\label{eq:nu-moments}
 \int_{\R}x^m\dd\nu(x)=0\qquad(m\ge1).
\end{equation}
Consequently $\nu-\tfrac12\delta_0$ is a nonzero finite signed measure whose
moments of every nonnegative integer order vanish.
\end{corollary}

\begin{proof}
The total variation and every polynomial moment are finite because $(a_k)$ is
Schwartz.  Odd moments vanish by symmetry; positive even moments vanish by
\cref{thm:weighted-prouhet}; and the mass is \eqref{eq:mass-half}.
\end{proof}

\begin{remark}
There is no contradiction with uniqueness results for positive moment problems:
\eqref{eq:nu-measure} is signed and has unbounded support.  It gives a concrete
moment-null object built from the exact spectrum of a positive compactly
supported density.
\end{remark}

\subsection{Certified finite truncations}

The identities are infinite, but the elementary product bound gives explicit
certificates for truncation experiments.

\begin{proposition}[Prouhet truncation bound]\label{prop:prouhet-tail}
Let $m\ge0$, $K\ge1$, and $J>2m+1$ be integers.  Put $p=J-2m$.  Then
\begin{align}
 \left|\sum_{k=K}^{\infty}a_ko_k^{2m}\right|
 &\le \frac{2^{J(J+1)/2}}{\pi^J}
 \sum_{k=K}^{\infty}(2k+1)^{-p}
 \label{eq:prouhet-tail-raw}\\
 &\le \frac{2^{J(J+1)/2}}{\pi^J}
 \left((2K+1)^{-p}
 +\frac{(2K+1)^{1-p}}{2(p-1)}\right).
 \label{eq:prouhet-tail-explicit}
\end{align}
For $m\ge1$, the same bound is an a posteriori certificate for the residual of
the first $K$ terms, because the complete sum is zero.
\end{proposition}

\begin{proof}
Apply \eqref{eq:Phi-power-bound-baseline} with $x=o_k/2$.  Since
\[
 |a_k|\le
 \frac{2^{J(J-1)/2}}{\pi^J(k+1/2)^J}
 =\frac{2^{J(J+1)/2}}{\pi^Jo_k^J},
\]
multiplication by $o_k^{2m}$ and summation give
\eqref{eq:prouhet-tail-raw}.  For the second line use the monotonicity of
$(2x+1)^{-p}$ and the integral estimate
\[
 \sum_{k=K}^{\infty}(2k+1)^{-p}
 \le (2K+1)^{-p}+\int_K^\infty(2x+1)^{-p}\dd x.
\]
\end{proof}

\begin{remark}
The coefficient in \eqref{eq:prouhet-tail-raw} records every factor of two from
the half-integer substitution $x=o_k/2$.  The sharper alias estimate in
\cref{sec:alias-error} works directly in the variable $x$ and therefore uses the
constant $C_K$ of \eqref{eq:selected-factor-bound} without this substitution.
\end{remark}

\section{Dyadic spectral moment filtration}\label{sec:moment-filtration}

\subsection{Every derivative as a half-integer spectral moment}

For $m\ge0$ define
\begin{equation}\label{eq:S-m-def}
 S_m(t):=\sum_{k=0}^{\infty}a_ko_k^m
 \cos\!\left(\pi o_kt+\frac{m\pi}{2}\right).
\end{equation}

\begin{theorem}[Spectral derivative identity]\label{thm:spectral-derivative}
For $-1\le t\le1$ and every $m\ge0$,
\begin{equation}\label{eq:spectral-derivative}
 \boxed{
 S_m(t)=\frac{\Up^{(m)}(t)}{\pi^m}
 =\frac{2^{\binom{m+1}{2}}}{\pi^m}
 \GF\bigl(2^m(t+1)\bigr).}
\end{equation}
The series is absolutely and uniformly convergent on $[-1,1]$.
\end{theorem}

\begin{proof}
Differentiate \eqref{eq:cosine-spectrum} $m$ times.  The derivative of
$\cos(\pi o_kt)$ is
$(\pi o_k)^m\cos(\pi o_kt+m\pi/2)$.  The second equality is the imported
global derivative identity \eqref{eq:global-derivatives}.
\end{proof}

\subsection{Exact termination at reduced dyadic points}

\begin{theorem}[Dyadic spectral filtration]\label{thm:dyadic-filtration}
Let $n\ge1$ and let $a$ be odd with $0<a<2^n$.  Put $t=a/2^n$.  Then
\begin{align}
 S_m(t)&=0 &&(m>n),
 \label{eq:filtration-above}\\
 S_n(t)&=-\eps_{(a-1)/2}
 \frac{2^{\binom{n+1}{2}}}{\pi^n},
 \label{eq:filtration-critical}\\
 S_m(t)&=\frac{2^{\binom{m+1}{2}}}{\pi^m}
 \GF\!\left(\frac{2^n+a}{2^{n-m}}\right)
 &&(0\le m<n).
 \label{eq:filtration-below}
\end{align}
Every value in \eqref{eq:filtration-below} is rational up to the displayed
power of $\pi$, and is given explicitly by \eqref{eq:baseline-qbinomial}.
In particular, the reduced dyadic depth $n$ is exactly the largest positive
order for which the spectral moment $S_m(t)$ can be nonzero.
\end{theorem}

\begin{proof}
By \cref{thm:spectral-derivative}, the relevant argument of $\GF$ is
\[
 2^m(t+1)=\frac{2^n+a}{2^{n-m}}.
\]
If $m>n$, this is an even integer, so \eqref{eq:global-integers} gives zero.
If $m=n$, the argument is the odd integer $2^n+a=2b+1$, with
\[
 b=2^{n-1}+\frac{a-1}{2}.
\]
The binary supports of the two summands are disjoint; hence
$\eps_b=-\eps_{(a-1)/2}$, proving \eqref{eq:filtration-critical}.  The remaining
case is simply the unreduced dyadic argument displayed in
\eqref{eq:filtration-below}, to which the global $q$-binomial formula applies.
The critical value is nonzero, while all larger orders vanish.
\end{proof}

\begin{corollary}[Spectral detection of the denominator]\label{cor:denominator-detection}
For an interior reduced dyadic point $t=a/2^n$ with odd $a$,
\begin{equation}\label{eq:denominator-detection}
 n=\max\{m\ge1:S_m(t)\ne0\}.
\end{equation}
Thus the denominator depth is encoded intrinsically in the half-integer
Fourier spectrum, without inspecting the binary expansion of $t$.
\end{corollary}

\begin{remark}
The underlying finite derivative jet at a dyadic point is baseline material in
the repository.  The new content is its exact re-expression as a hierarchy of
weighted half-integer spectral moments and the resulting denominator detector.
\end{remark}

\section{An exact beyond-all-orders spectral remainder}\label{sec:beyond-all-orders}

\subsection{Cancellation of every cosine Taylor coefficient}

For $M\ge0$ define
\begin{equation}\label{eq:rho-M}
 \rho_M(y):=1-\cos y-
 \sum_{m=1}^{M}(-1)^{m+1}\frac{y^{2m}}{(2m)!},
\end{equation}
with the sum empty when $M=0$.

\begin{theorem}[Exact all-order spectral remainder]\label{thm:spectral-remainder}
For every $x\in[0,1]$ and every integer $M\ge0$,
\begin{equation}\label{eq:F-basic-spectral}
 F(x)=\sum_{k=0}^{\infty}a_k
 \bigl(1-\cos(\pi o_kx)\bigr),
\end{equation}
and, more strongly,
\begin{equation}\label{eq:F-rho-spectral}
 \boxed{
 F(x)=\sum_{k=0}^{\infty}\eps_kw_k\,
 \rho_M(\pi o_kx).}
\end{equation}
The series in \eqref{eq:F-rho-spectral} converges absolutely for every fixed
$M$ and uniformly for $x\in[0,1]$.
\end{theorem}

\begin{proof}
For $0\le x\le1$, \eqref{eq:Up-F-relations} gives
$F(x)=1-\Up(x)$.  Subtract \eqref{eq:cosine-spectrum} from $1$ and use
$\sum_ka_k=1/2$ to obtain \eqref{eq:F-basic-spectral}.  Subtracting the first
$M$ nonconstant Taylor terms of $1-\cos y$ changes the right side by
\[
 \sum_{m=1}^{M}\frac{(-1)^{m+1}(\pi x)^{2m}}{(2m)!}
 \sum_{k=0}^{\infty}a_ko_k^{2m},
\]
which is zero by \cref{thm:weighted-prouhet}.  Rapid decay of $a_k$ dominates
the polynomial growth of $\rho_M(\pi o_kx)$ uniformly on $[0,1]$.
\end{proof}

\begin{remark}[A spectral explanation of flatness]
For a fixed individual frequency, $1-\cos(\pi o_kx)$ has an ordinary even
Taylor series at zero.  Formula \eqref{eq:F-rho-spectral} says that after the
frequencies are summed with their exact Thue--Morse sinc weights, \emph{every}
algebraic Taylor coefficient disappears.  The endpoint value survives only as
a collective remainder beyond all finite algebraic orders.  This is stronger
than checking $F^{(m)}(0)=0$ one derivative at a time: the same exact identity
holds after subtraction at any prescribed order $M$.
\end{remark}

\subsection{The direct Lambert--\texorpdfstring{$W$}{W} bridge}

\begin{corollary}[Thue--Morse--sinc--Lambert bridge]\label{cor:Lambert-spectral-bridge}
Let $M\ge0$ be fixed, $L=\log2$, and
$W=W_{-1}(-Lx)$.  As $x\downarrow0$,
\begin{equation}\label{eq:Lambert-spectral-bridge}
\boxed{
\begin{aligned}
 \sum_{k=0}^{\infty}\eps_kw_k\rho_M(\pi o_kx)
 ={}&\frac{\e^{\Psi(-W/L)}}{2^{7/12}\sqrt{\pi x}}
 \exp\!\left(\frac{\mathfrak g-W-\tfrac12W^2}{L}\right)\\
 &\times\left(1+\bigO\!\left(\lambda(x)^{-1}\right)\right).
\end{aligned}}
\end{equation}
The left side is independent of $M$, although its individual summands are not.
\end{corollary}

\begin{proof}
The left side equals $F(x)$ exactly by
\cref{thm:spectral-remainder}.  Apply the imported endpoint theorem
\eqref{eq:Lambert-baseline}.
\end{proof}

\begin{remark}[What has and has not been proved]
Equation \eqref{eq:Lambert-spectral-bridge} is a rigorous bridge, but it does not
yet identify a sharply localized set of indices $k$ that carries the asymptotic.
The natural lower-Lambert scale $r(x)=\lambda(x)/x$ strongly suggests such a
localization problem.  Establishing it requires estimates that preserve the
massive cancellations in \eqref{eq:F-rho-spectral}; termwise absolute values
would destroy the phenomenon.  We formulate this as an open problem in
\cref{sec:open}.
\end{remark}

\section{The exact dyadic alias transform}\label{sec:dct}

\subsection{Folded coefficients}

Fix $n\ge1$ and put $N=2^n$.  For every integer $r$, define the
length-$2N$ folded Fourier coefficient
\begin{equation}\label{eq:A-n-r-def}
 A_{n,r}:=\frac1N\sum_{j=0}^{2N-1}
 P\!\left(\frac jN\right)\e^{-\pi\ii rj/N}.
\end{equation}
It depends only on $r\bmod 2N$ and is real because $P$ is real and even.  The
nontrivial coefficients are represented by odd $r$ with $1\le r<N$.

\begin{theorem}[Exact half-integer alias identity]\label{thm:alias-identity}
For every admissible $n,r$,
\begin{equation}\label{eq:alias-identity}
 \boxed{
 A_{n,r}=\sum_{q\in\Z}\PhiU\!\left(qN+\frac r2\right).}
\end{equation}
The series converges absolutely.  Moreover,
\begin{align}
 A_{n,r}
 &=\frac1N\left[1+2\sum_{j=1}^{N-1}
 \Up\!\left(\frac jN\right)
 \cos\!\left(\frac{\pi rj}{N}\right)\right]
 \label{eq:A-Up-DCT}\\
 &=\frac1N\left[1-2\sum_{j=1}^{N-1}
 F\!\left(\frac jN\right)
 \cos\!\left(\frac{\pi rj}{N}\right)\right].
 \label{eq:A-F-DCT}
\end{align}
Thus a finite discrete cosine transform of exact dyadic Fabius values equals a
complete folded alias class of the infinite sinc product.
\end{theorem}

\begin{proof}
Write the complex Fourier expansion of $P$ as
\[
 P(t)=\sum_{m\in\Z}c_m\e^{\pi\ii mt},
 \qquad c_m=\frac12\PhiU(m/2).
\]
Substitute it in \eqref{eq:A-n-r-def}.  Absolute convergence permits
interchange, and the finite character sum is $2N$ when $m\equiv r\pmod{2N}$
and zero otherwise.  Therefore
\[
 A_{n,r}=2\sum_{m\equiv r\ (2N)}c_m
 =\sum_{q\in\Z}\PhiU\!\left(\frac{r+2Nq}{2}\right),
\]
which is \eqref{eq:alias-identity}.  Absolute convergence follows from rapid
decay.

Pair the terms $j$ and $2N-j$ in \eqref{eq:A-n-r-def}.  The endpoint
$j=N$ contributes zero because $P(1)=\Up(1)=0$, while $P(0)=1$; this gives
\eqref{eq:A-Up-DCT}.  On $0\le j\le N$,
$\Up(j/N)=1-F(j/N)$.  Since $r$ is odd,
\[
 \sum_{j=1}^{N-1}\cos(\pi rj/N)=0,
\]
which yields \eqref{eq:A-F-DCT}.
\end{proof}

\subsection{Exact inverse transform}

\begin{theorem}[Dyadic spectral reconstruction]\label{thm:inverse-DCT}
For every integer $0\le j\le2N-1$,
\begin{equation}\label{eq:inverse-DCT}
 \boxed{
 P\!\left(\frac jN\right)
 =\frac12+
 \sum_{\substack{1\le r<N\\r\ \mathrm{odd}}}
 A_{n,r}\cos\!\left(\frac{\pi rj}{N}\right).}
\end{equation}
In particular, all dyadic values of $\Up$ at denominator $N$ are reconstructed
from the $N/2$ folded half-integer classes.
\end{theorem}

\begin{proof}
The normalization in \eqref{eq:A-n-r-def} is twice the standard length-$2N$
DFT normalization.  Its inverse is therefore
\[
 P(j/N)=\frac12\sum_{r=0}^{2N-1}A_{n,r}\e^{\pi\ii rj/N}.
\]
For $r=0$, the alias sum contains $\PhiU(0)=1$ and only zero-valued nonzero
integer samples, so $A_{n,0}=1$.  For nonzero even $r$, every alias argument is
a nonzero integer and $A_{n,r}=0$.  Finally
$A_{n,2N-r}=A_{n,r}$.  Pairing the odd frequencies gives
\eqref{eq:inverse-DCT}.
\end{proof}

\begin{corollary}[Exact discrete energy identity]\label{cor:discrete-energy}
One has
\begin{equation}\label{eq:discrete-energy}
 \boxed{
 \sum_{\substack{1\le r<N\\r\ \mathrm{odd}}}A_{n,r}^2
 =\frac1N\left[1+2\sum_{j=1}^{N-1}
 \Up\!\left(\frac jN\right)^2\right]-\frac12.}
\end{equation}
The right side is rational.
\end{corollary}

\begin{proof}
Apply finite Parseval to \eqref{eq:inverse-DCT}.  The constant mode contributes
$1/4$, and each nonconstant cosine has mean square $1/2$ on the length-$2N$
grid.  Use $P(0)=1$, $P(1)=0$, and symmetry of the remaining samples.
\end{proof}

\begin{remark}[Continuous limit]
As $n\to\infty$, the right side of \eqref{eq:discrete-energy} tends
\[
 2\int_0^1\Up(t)^2\dd t-\frac12
 =\sum_{k=0}^{\infty}a_k^2,
\]
which is also ordinary Fourier Parseval for \eqref{eq:cosine-spectrum}.  The
finite identity is stronger: before taking a limit it expresses the complete
Galois-square trace of the folded spectrum using rational dyadic samples.
\end{remark}

\section{Certified superexponential recovery}\label{sec:alias-error}

\subsection{A uniform product estimate}

\begin{lemma}[Selected-factor bound]\label{lem:selected-factor-bound}
For every integer $K\ge1$ and real $x\ne0$,
\begin{equation}\label{eq:selected-factor-bound}
 |\PhiU(x)|\le C_K|x|^{-K},
 \qquad
 C_K:=\frac{2^{K(K-1)/2}}{\pi^K}.
\end{equation}
\end{lemma}

\begin{proof}
For $0\le j<K$,
\[
 \left|\sinc(x/2^j)\right|
 \le\frac{2^j}{\pi|x|}.
\]
Multiply these $K$ inequalities and bound every remaining sinc factor by one.
\end{proof}

\begin{theorem}[Uniform aliasing error]\label{thm:alias-error}
Let $K>1$ be an integer.  For $N=2^n$ and odd $1\le r<N$,
\begin{equation}\label{eq:alias-error-general}
 \boxed{
 |A_{n,r}-\PhiU(r/2)|
 \le 2C_KN^{-K}(2^K-1)\zeta(K).}
\end{equation}
Choosing $K=n\ge2$ gives the explicit uniform estimate
\begin{equation}\label{eq:alias-error-n}
 \boxed{
 |A_{n,r}-\PhiU(r/2)|
 \le E_n:=
 \frac{2(1-2^{-n})\zeta(n)}{\pi^n}
 2^{-n(n-1)/2}.}
\end{equation}
In particular $E_n=2^{-n^2/2+\bigO(n\log n)}$, uniformly over all odd
$r<N$.
\end{theorem}

\begin{proof}
Remove the $q=0$ term from \eqref{eq:alias-identity}.  For $q\ge1$ and
$0<r<N$,
\[
 |qN\pm r/2|\ge(q-1/2)N.
\]
Hence \cref{lem:selected-factor-bound} gives
\begin{align*}
 |A_{n,r}-\PhiU(r/2)|
 &\le 2C_KN^{-K}\sum_{q=1}^{\infty}(q-1/2)^{-K}\\
 &=2C_KN^{-K}(2^K-1)\zeta(K),
\end{align*}
using
$\sum_{q\ge1}(q-1/2)^{-K}=(2^K-1)\zeta(K)$.
Substitute $N=2^n$ and $K=n$; collecting powers of two gives
\eqref{eq:alias-error-n}.
\end{proof}

\begin{remark}[Why the rate is so fast]
At level $n$, the unwanted aliases begin a distance comparable with $2^n$
from the target.  The same sinc product permits an $n$-factor power bound, so
the alias tail behaves like $(2^n)^{-n}$ multiplied by the factor
$2^{n(n-1)/2}$ accumulated from dyadic scaling.  The net exponent is
$-n(n-1)/2$.
\end{remark}

\subsection{Numerical convergence for the leading amplitude}

The distinguished coefficient $A_{n,1}$ converges to
\begin{equation}\label{eq:h-value}
 h:=\PhiU(1/2)
 =0.553771275888810095344217416344548881284886642554\ldots.
\end{equation}
The finite values in \cref{tab:A-convergence} were computed from exact rational
$F(j/2^n)$ values, with only the final cyclotomic evaluation performed at high
precision.

\begin{table}[ht]
\centering\small
\begin{tabular}{@{}r r r r@{}}
\toprule
$n$ & $A_{n,1}$ & actual error & rigorous $E_n$ \\
\midrule
2 & 0.5544487530108746285614747 & $6.7748\,10^{-4}$ & $1.2500\,10^{-1}$ \\
3 & 0.5537613501982070710846451 & $9.9257\,10^{-6}$ & $8.4805\,10^{-3}$ \\
4 & 0.5537712795003417162450708 & $3.6115\,10^{-9}$ & $3.2552\,10^{-4}$ \\
5 & 0.5537712758798224271695171 & $8.9877\,10^{-12}$ & $6.4112\,10^{-6}$ \\
6 & 0.5537712758888101658815125 & $7.0537\,10^{-17}$ & $6.3578\,10^{-8}$ \\
7 & 0.5537712758888100953109457 & $3.3272\,10^{-20}$ & $3.1590\,10^{-10}$ \\
8 & 0.5537712758888100953442174 & $5.0949\,10^{-27}$ & $7.8534\,10^{-13}$ \\
9 & 0.5537712758888100953442174 & $4.8386\,10^{-31}$ & $9.7639\,10^{-16}$ \\
10& 0.5537712758888100953442174 & $1.3837\,10^{-39}$ & $6.0700\,10^{-19}$ \\
\bottomrule
\end{tabular}
\caption{Superexponential recovery of the first half-integer amplitude.  The
bound is deliberately elementary and uniform in $r$, so it is not expected to
track the much smaller observed error closely.}
\label{tab:A-convergence}
\end{table}

\section{Finite \texorpdfstring{$q$}{q}-binomial--cyclotomic approximants}\label{sec:q-approximants}

\subsection{Explicit substitution}

The alias transform is finite, while every dyadic value in it has the finite
Gaussian-binomial form \eqref{eq:baseline-qbinomial}.  Their composition is the
promised bridge between the two coefficient webs.

\begin{theorem}[$q$-binomial half-integer approximant]\label{thm:q-approximant}
Let $N=2^n$ and let $r$ be odd with $1\le r<N$.  Then
\begin{equation}\label{eq:q-approximant}
\boxed{
\begin{aligned}
 A_{n,r}={}&\frac1N-
 \frac{2}{N\,2^{n^2}(1/2;1/2)_n}
 \sum_{j=1}^{N-1}\cos\!\left(\frac{\pi rj}{N}\right)\\
 &\quad\times\sum_{k=0}^{n}
 \frac{\qbinom{n}{k}{1/2}}
 {2^{k(k-1)}(n+k)!}
 \sum_{s=0}^{j2^k-1}\eps_s(s-j2^k)^{n+k}.
\end{aligned}}
\end{equation}
For every fixed odd $r$, the right side converges to $\PhiU(r/2)$ as
$n\to\infty$ through levels with $r<2^n$, and the error is at most $E_n$.
\end{theorem}

\begin{proof}
Insert \eqref{eq:baseline-qbinomial} with $m=j$ and $c=0$ into
\eqref{eq:A-F-DCT}.  The convergence statement is
\cref{thm:alias-error}.
\end{proof}

\begin{remark}[Three structures in one finite expression]
Formula \eqref{eq:q-approximant} simultaneously contains
\begin{enumerate}
 \item Gaussian binomial coefficients at $q=1/2$;
 \item finite Thue--Morse power sums;
 \item real $2^{n+1}$-st cyclotomic numbers
 $\cos(\pi rj/2^n)$.
\end{enumerate}
No infinite product or numerical quadrature is needed to construct the
approximant.  The infinite sinc product enters only as the certified limit.
\end{remark}

\subsection{A finite algebraic tower}

Let
\begin{equation}\label{eq:cyclotomic-field}
 \zeta_{2N}=\e^{\pi\ii/N},
 \qquad
 K_n:=\Q(\zeta_{2N}+\zeta_{2N}^{-1}).
\end{equation}
Since every $F(j/N)$ is rational,
\begin{equation}\label{eq:A-in-K}
 A_{n,r}\in K_n.
\end{equation}
The fields form the nested real dyadic cyclotomic tower
$K_1\subset K_2\subset\cdots$.  The approximants therefore approach the
transcendence-unknown numbers $\PhiU(r/2)$ through an explicit tower of exact
algebraic numbers with a superexponential analytic certificate.

\begin{remark}
The phrase ``transcendence-unknown'' is intentional.  No irrationality or
transcendence assertion for $h=\PhiU(1/2)$ is proved here or in the reviewed
corpus.  The algebraic approximants are useful precisely because they separate
exact finite arithmetic from the unresolved nature of the limiting constants.
\end{remark}

\section{Cyclotomic orbit polynomials and rational traces}\label{sec:cyclotomic}

\subsection{The Galois orbit}

The real cyclotomic Galois group is
\begin{equation}\label{eq:G-n}
 G_n\simeq(\Z/2N\Z)^{\times}/\{\pm1\}.
\end{equation}
Its elements may be represented by odd integers $u$ modulo $2N$, with $u$ and
$-u$ identified.  The representatives $1\le r<N$, $r$ odd, form a complete
set.

\begin{proposition}[Galois action]\label{prop:Galois-action}
If $\sigma_u(\zeta_{2N})=\zeta_{2N}^u$, then
\begin{equation}\label{eq:Galois-action}
 \sigma_u(A_{n,r})=A_{n,r'},
\end{equation}
where $r'$ is the unique odd representative in $[1,N)$ congruent to
$\pm ur\pmod{2N}$.
\end{proposition}

\begin{proof}
Apply $\sigma_u$ to \eqref{eq:A-F-DCT}.  It fixes the rational Fabius values and
sends
$\cos(\pi rj/N)$ to $\cos(\pi urj/N)$.  Reduction modulo $2N$ and evenness of
cosine give the stated representative.
\end{proof}

\begin{definition}[Orbit polynomial and power traces]\label{def:orbit-poly}
Define
\begin{equation}\label{eq:orbit-poly}
 \mathcal P_n(X):=
 \prod_{\substack{1\le r<N\\r\ \mathrm{odd}}}(X-A_{n,r})
\end{equation}
and, for $m\ge1$,
\begin{equation}\label{eq:p-m-n}
 p_m(n):=\sum_{\substack{1\le r<N\\r\ \mathrm{odd}}}A_{n,r}^m.
\end{equation}
The polynomial is called an orbit polynomial rather than a minimal polynomial:
if the stabilizer of $A_{n,1}$ is nontrivial, repeated conjugates may occur.
\end{definition}

\begin{theorem}[Rational orbit polynomial]\label{thm:rational-orbit-poly}
For every $n\ge1$,
\begin{equation}\label{eq:orbit-poly-rational}
 \mathcal P_n(X)\in\Q[X],
 \qquad
 p_m(n)\in\Q\quad(m\ge1).
\end{equation}
The coefficients of $\mathcal P_n$ are recovered from the $p_m(n)$ by Newton's
identities.
\end{theorem}

\begin{proof}
The multiset of roots is invariant under the full Galois group by
\cref{prop:Galois-action}; hence every elementary symmetric function is fixed
by $G_n$ and lies in $\Q$.  The same argument applies to power sums.  Newton's
identities are purely rational algebra.
\end{proof}

\subsection{Ramanujan-filtered trace formula}

Put
\begin{equation}\label{eq:P-j-sample}
 P_j:=P(j/N),
 \qquad 0\le j<2N.
\end{equation}
Every $P_j$ is rational, since it is a dyadic value of $\Up$.

\begin{theorem}[Exact rational power traces]\label{thm:Ramanujan-traces}
For every $m\ge1$,
\begin{equation}\label{eq:Ramanujan-traces}
\boxed{
 p_m(n)=\frac12N^{1-m}
 \sum_{\substack{0\le j_1,\ldots,j_m<2N\\
                  N\mid j_1+\cdots+j_m}}
 (-1)^{(j_1+\cdots+j_m)/N}
 \prod_{\ell=1}^{m}P_{j_\ell}.}
\end{equation}
In particular,
\begin{equation}\label{eq:universal-trace}
 \boxed{p_1(n)=\frac12\qquad(n\ge1).}
\end{equation}
For $m=2$, \eqref{eq:Ramanujan-traces} reduces to the discrete energy identity
\eqref{eq:discrete-energy}.
\end{theorem}

\begin{proof}
Expand $A_{n,r}^m$ from \eqref{eq:A-n-r-def} and sum over the $N/2$ odd
representatives.  Since $A_{n,r}=A_{n,2N-r}$, the representative sum is one
half of the sum over all odd residues modulo $2N$.  Thus
\begin{align*}
 p_m(n)
 &=\frac{1}{2N^m}
 \sum_{0\le j_1,\ldots,j_m<2N}
 \left(\prod_{\ell=1}^mP_{j_\ell}\right)
 c_{2N}(j_1+\cdots+j_m),
\end{align*}
where
\[
 c_{2N}(S)=\sum_{\substack{r\bmod 2N\\r\ \mathrm{odd}}}
 \e^{-\pi\ii rS/N}
\]
is the Ramanujan sum for the power-of-two modulus $2N$.  Its elementary value is
\begin{equation}\label{eq:Ramanujan-power-two}
 c_{2N}(S)=
 \begin{cases}
  N(-1)^{S/N},&N\mid S,\\
  0,&N\nmid S.
 \end{cases}
\end{equation}
Substitution gives \eqref{eq:Ramanujan-traces}.  For $m=1$, the only indices
with $N\mid j$ are $j=0,N$, so
$p_1(n)=\tfrac12(P_0-P_N)=1/2$.
\end{proof}

\begin{remark}[A trace that does not disperse]
The orbit degree is $N/2\to\infty$, yet the unnormalized first trace remains
exactly $1/2$.  Most conjugates become tiny; a sparse collection of low-frequency
conjugates carries the persistent trace.  This is quantified below.
\end{remark}

\subsection{Orbit polynomials at the first levels}

At $n=1$ there is one coefficient, $A_{1,1}=1/2$, so
\begin{equation}\label{eq:P1-poly}
 \mathcal P_1(X)=X-\frac12.
\end{equation}
At $n=2$,
\begin{equation}\label{eq:A2-values}
 A_{2,1}=\frac14+\frac{31\sqrt2}{144},
 \qquad
 A_{2,3}=\frac14-\frac{31\sqrt2}{144},
\end{equation}
and
\begin{equation}\label{eq:P2-poly}
 \boxed{
 \mathcal P_2(X)=X^2-\frac12X-\frac{313}{10368}.}
\end{equation}
At $n=3$, direct use of \eqref{eq:Ramanujan-traces} and Newton identities gives
\begin{equation}\label{eq:P3-poly}
\boxed{
\begin{aligned}
 \mathcal P_3(X)={}&X^4-\frac12X^3
 -\frac{1627}{55296}X^2
 -\frac{4603}{23887872}X
 +\frac{32257}{220150628352}.
\end{aligned}}
\end{equation}
The corresponding rational power traces are
\begin{align}
 p_1(3)&=\frac12,
 &p_2(3)&=\frac{8539}{27648},
 \label{eq:p3-12}\\
 p_3(3)&=\frac{1351363}{7962624},
 &p_4(3)&=\frac{5175727235}{55037657088}.
 \label{eq:p3-34}
\end{align}
The four conjugates are approximately
\begin{equation}\label{eq:n3-conjugates}
\begin{aligned}
 &\phantom{-}0.553761350198207071,
 &&-0.046098883415321902,\\
 &-0.008349869595552727,
 &&\phantom{-}0.000687402812667557.
\end{aligned}
\end{equation}
They already resemble the first four half-integer amplitudes, but each includes
all aliases in its congruence class.

At $n=4$ the eight conjugates are approximately
\begin{align*}
 &\phantom{-}0.553771279500341716,
 &&-0.046202534866953151,\\
 &-0.008657957107189789,
 &&\phantom{-}0.001046027211381176,\\
 &-0.000358624398713618,
 &&\phantom{-}0.000308087511637062,\\
 &\phantom{-}0.000103651451631249,
 &&-0.000009929302134645.
\end{align*}
Their sum is exactly $1/2$.

\subsection{Rational traces converging to nonlinear sinc moments}

Define the nonlinear half-integer spectral sums
\begin{equation}\label{eq:spectral-power-sums}
 \mathfrak s_m:=\sum_{k=0}^{\infty}a_k^m
 \qquad(m\ge1).
\end{equation}
They converge absolutely.  The case $m=1$ is $\mathfrak s_1=1/2$; for $m=2$,
\begin{equation}\label{eq:s2-L2}
 \mathfrak s_2=2\int_0^1\Up(t)^2\dd t-\frac12.
\end{equation}

\begin{theorem}[Cyclotomic trace limits]\label{thm:trace-limits}
For every fixed integer $m\ge1$,
\begin{equation}\label{eq:trace-limit}
 \boxed{
 \lim_{n\to\infty}p_m(n)=\mathfrak s_m.}
\end{equation}
For $m=1$ the equality is exact at every level.  For fixed $m\ge2$, the error
splits into
\begin{equation}\label{eq:trace-error-split}
 |p_m(n)-\mathfrak s_m|
 \le \frac{Nm}{2}(1+E_n)^{m-1}E_n
 +\sum_{k\ge N/2}|a_k|^m,
\end{equation}
where both terms decay faster than every exponential in $n$ after applying an
$n$-factor sinc bound to the tail.
\end{theorem}

\begin{proof}
Index the admissible residues by $r=2k+1$, $0\le k<N/2$.  Then
\cref{thm:alias-error} gives
$|A_{n,2k+1}-a_k|\le E_n$.  Since $|a_k|\le1$ and $E_n\le1/8$ for $n\ge2$,
the mean-value inequality for $x\mapsto x^m$ gives
\[
 |A_{n,2k+1}^m-a_k^m|
 \le m(1+E_n)^{m-1}E_n.
\]
Summing over $k<N/2$ and adding the omitted tail proves
\eqref{eq:trace-error-split}.  The first term is superexponential because
$NE_n=2^{-n^2/2+\bigO(n)}$.  For the tail, use
\cref{lem:selected-factor-bound} with $K=n$:
\[
 |a_k|\le C_n(k+1/2)^{-n}.
\]
The resulting integral bound is
\[
 \bigO\!\left(C_n^m(N/2)^{1-nm}\right)
 =2^{-mn^2/2+\bigO(n)},
\]
where the implied constant may depend on the fixed exponent $m$.
\end{proof}

\begin{corollary}[Algebraic approximation to spectral norms]\label{cor:algebraic-spectral-norms}
Every $\mathfrak s_m$ is a limit of explicitly computable rational numbers
$p_m(n)$ whose defining algebraic conjugates are given by the finite
$q$-binomial formula \eqref{eq:q-approximant}.  In particular,
$2\int_0^1\Up^2-1/2$ admits a superexponentially convergent rational sequence
coming from cyclotomic Galois square traces.
\end{corollary}

\subsection{Collapse of the normalized Galois orbit}

Let
\begin{equation}\label{eq:empirical-measure}
 \mu_n:=\frac2N
 \sum_{\substack{1\le r<N\\r\ \mathrm{odd}}}\delta_{A_{n,r}}.
\end{equation}
This is a probability measure on $\R$.

\begin{theorem}[Galois-orbit collapse]\label{thm:orbit-collapse}
The measures $\mu_n$ converge weakly to the point mass $\delta_0$.  More
quantitatively,
\begin{equation}\label{eq:orbit-second-moment}
 \int x^2\dd\mu_n(x)=\frac2N p_2(n)\le\frac3N.
\end{equation}
Nevertheless the unnormalized trace remains $p_1(n)=1/2$.
\end{theorem}

\begin{proof}
From \eqref{eq:discrete-energy} and $0\le\Up\le1$,
\[
 p_2(n)\le\frac1N(1+2(N-1))-\frac12<\frac32.
\]
This proves \eqref{eq:orbit-second-moment}.  Chebyshev's inequality shows that
the proportion of conjugates outside any fixed neighborhood of zero tends to
zero, hence $\mu_n\Rightarrow\delta_0$.
\end{proof}

\begin{remark}
This is an unusual arithmetic picture: the degree of the orbit grows, the
normalized orbit collapses to zero, but its unnormalized first trace remains
fixed.  The low-index half-integer amplitudes survive as a vanishing proportion
of the conjugates and carry the entire trace in the limit.
\end{remark}

\section{A finite Thue--Morse sign carrier for every alias class}\label{sec:sign-carrier}

\subsection{Exact factorization of a folded class}

Let $N=2^n$, let $r$ be odd with $1\le r<N$, and put
\begin{equation}\label{eq:a-shift}
 \alpha:=\frac{r}{2N}\in(0,1/2).
\end{equation}
Define the finite prefactor
\begin{equation}\label{eq:Gamma-n-r}
 \Gamma_{n,r}:=
 \frac{\displaystyle\prod_{h=1}^{n}\sin(\pi r/2^h)}
      {\pi^n2^{n(n+1)/2}}.
\end{equation}

\begin{theorem}[Alias-class factorization]\label{thm:alias-factorization}
For every $q\in\Z$,
\begin{equation}\label{eq:single-alias-factorization}
 \PhiU\!\left(qN+\frac r2\right)
 =\Gamma_{n,r}\,
 \frac{\PhiU(q+\alpha)}{(q+\alpha)^n}.
\end{equation}
Consequently
\begin{equation}\label{eq:A-shifted-sum}
 \boxed{
 A_{n,r}=\Gamma_{n,r}
 \sum_{q\in\Z}\frac{\PhiU(q+\alpha)}{(q+\alpha)^n}.}
\end{equation}
The shifted sum converges absolutely.
\end{theorem}

\begin{proof}
Split the product \eqref{eq:Phi-def} after the first $n$ factors.  Since
$qN+r/2=N(q+\alpha)$, the tail is
\[
 \prod_{j=n}^{\infty}
 \sinc\!\left(\frac{qN+r/2}{2^j}\right)
 =\PhiU(q+\alpha).
\]
For $0\le j<n$,
\[
 \sin\!\left(\pi\frac{qN+r/2}{2^j}\right)
 =\sin\!\left(\pi q2^{n-j}+\frac{\pi r}{2^{j+1}}\right)
 =\sin\!\left(\frac{\pi r}{2^{j+1}}\right),
\]
because $q2^{n-j}$ is even.  The denominators of the first $n$ sinc factors
contribute
\[
 \prod_{j=0}^{n-1}\frac{2^j}{\pi N(q+\alpha)}
 =\frac{1}{\pi^n2^{n(n+1)/2}(q+\alpha)^n}.
\]
Multiplication proves \eqref{eq:single-alias-factorization}; summing over $q$
and using \eqref{eq:alias-identity} gives \eqref{eq:A-shifted-sum}.
\end{proof}

\subsection{The finite Thue--Morse polynomial}

Let
\begin{equation}\label{eq:T-n-polynomial}
 T_n(z):=\prod_{\ell=0}^{n-1}(1-z^{2^\ell})
 =\sum_{m=0}^{2^n-1}\eps_mz^m.
\end{equation}

\begin{proposition}[Root-of-unity magnitude and sign]\label{prop:TM-sign-carrier}
For odd $1\le r<N$,
\begin{align}
 \left|T_n\!\left(\e^{2\pi\ii r/N}\right)\right|
 &=2^n\left|\prod_{h=1}^{n}\sin(\pi r/2^h)\right|,
 \label{eq:T-root-magnitude}\\
 \operatorname{sgn}\Gamma_{n,r}
 &=(-1)^{\sum_{h=1}^{n}\lfloor r/2^h\rfloor}
 =\eps_{(r-1)/2}.
 \label{eq:Gamma-sign}
\end{align}
Thus the finite prefactor in \eqref{eq:A-shifted-sum} carries exactly the
Thue--Morse sign attached to the target amplitude $\PhiU(r/2)$.
\end{proposition}

\begin{proof}
For the magnitude, use $|1-\e^{2\pi\ii x}|=2|\sin\pi x|$ in each factor of
\eqref{eq:T-n-polynomial} and reverse the order of scales.  For the sign,
$\sin(\pi x)$ has sign $(-1)^{\lfloor x\rfloor}$ away from integers.  Since
$r<N$, all floors beyond $h=n$ vanish, and
\[
 \sum_{h=1}^{n}\left\lfloor\frac r{2^h}\right\rfloor
 =r-s_2(r).
\]
Write $r=2c+1$.  Then $s_2(r)=s_2(c)+1$, so
$r-s_2(r)\equiv s_2(c)\pmod2$.  Hence the sign is $\eps_c$ with
$c=(r-1)/2$.
\end{proof}

\begin{remark}
The earlier arithmetic-ray report uses root-of-unity values of finite
Thue--Morse polynomials to control products along rational doubling orbits.  The
mechanism here is different: a finite root-of-unity product is extracted from a
\emph{Poisson alias class} of the sinc product.  The same automatic sign word
therefore governs both real-axis intervals and finite dyadic spectral folding.
\end{remark}

\subsection{A positivity conjecture}

Computations through $n=7$ show
\begin{equation}\label{eq:observed-A-sign}
 \operatorname{sgn}A_{n,r}=\eps_{(r-1)/2}
\end{equation}
for every admissible $r$.  By \eqref{eq:Gamma-sign}, this is equivalent to
positivity of the residual shifted sum.

\begin{conjecture}[Shifted spectral positivity]\label{conj:shifted-positivity}
For every $n\ge1$ and odd $1\le r<2^n$, with
$\alpha=r/2^{n+1}$,
\begin{equation}\label{eq:shifted-positivity}
 \boxed{
 \sum_{q\in\Z}\frac{\PhiU(q+\alpha)}{(q+\alpha)^n}>0.}
\end{equation}
Equivalently, every folded coefficient has the exact Thue--Morse sign
\eqref{eq:observed-A-sign}.

\smallskip\noindent\openstatus
\end{conjecture}

\begin{remark}[Why the conjecture is nontrivial]
The summands in \eqref{eq:shifted-positivity} have mixed signs.  Neither the
positivity of $\Up$ nor the elementary alias error bound proves the claim,
because for residues near $N$ the target amplitude can be much smaller than the
uniform alias bound.  A proof may require a quasi-periodic Green-kernel
representation or a total-positivity property of the dyadic sample sequence.
\end{remark}

\section{Certified arithmetic of integer-zero jets}\label{sec:zero-jets}

\subsection{The imported local formula}

Let $m\ne0$ be an integer and write
\begin{equation}\label{eq:m-v-q}
 m=2^vq,
 \qquad q\ \text{odd},
 \qquad d=v+1.
\end{equation}
The spectral determinant report proves the exact leading-jet identity
\begin{equation}\label{eq:leading-jet-baseline}
 \PhiU^{(d)}(m)=-\frac{d!}{m^d}\PhiU(q/2).
\end{equation}
It also gives every higher coefficient, but the leading term already shows why
the half-integer amplitudes are arithmetically central: every nonzero integer
zero reduces to one odd half-integer value.

\subsection{Finite exact approximants with a rigorous error}

\begin{theorem}[$q$-binomial approximation of every leading zero jet]\label{thm:jet-approximant}
Let $m=2^vq\ne0$ with $q$ odd and $d=v+1$.  Choose $n\ge2$ with
$2^n>|q|$, and put $r=|q|$.  Define
\begin{equation}\label{eq:J-m-n}
 J_{m,n}:=-\frac{d!}{m^d}A_{n,r}.
\end{equation}
Then $J_{m,n}$ is an explicitly computable element of the real dyadic
cyclotomic field $K_n$, given by the finite $q$-binomial expression
\eqref{eq:q-approximant}, and
\begin{equation}\label{eq:jet-error}
 \boxed{
 |\PhiU^{(d)}(m)-J_{m,n}|
 \le\frac{d!}{|m|^d}E_n.}
\end{equation}
Thus every leading integer-zero jet has a superexponentially convergent sequence
of exact algebraic approximants.
\end{theorem}

\begin{proof}
Evenness gives $\PhiU(q/2)=\PhiU(r/2)$.  Combine the imported identity
\eqref{eq:leading-jet-baseline} with \cref{thm:alias-error} and multiply the
error by $d!/|m|^d$.
\end{proof}

\begin{example}[The zero at $m=12$]\label{ex:jet-12}
Here $12=2^2\cdot3$, so $d=3$ and
\[
 \PhiU^{(3)}(12)=-\frac{6}{12^3}\PhiU(3/2)
 =\frac{-\PhiU(3/2)}{288}.
\]
Using the level-$4$ folded coefficient
$A_{4,3}=-0.04620253486695315055\ldots$ gives
\[
 J_{12,4}=0.000160425468288031773\ldots,
\]
while the exact target is
\[
 \PhiU^{(3)}(12)=0.000160424649771898309\ldots.
\]
The rigorous error follows immediately from \eqref{eq:jet-error}.  Increasing
$n$ replaces the displayed algebraic number by a much sharper one without
changing the proof architecture.
\end{example}

\begin{warning}
The final decimal in the target value in \cref{ex:jet-12} is included only as a
numerical orientation.  The certified content is the exact formula and the
bound \eqref{eq:jet-error}; the report's reproducibility code computes as many
digits as desired.
\end{warning}

\section{Computational verification}\label{sec:verification}

\subsection{Independent checks performed}

The numerical work was designed to check identities at several logically
independent interfaces rather than merely reproduce one formula in two notations.
The following were evaluated using exact rational arithmetic whenever possible
and 80--100 decimal digits for the final transcendental operations.
\begin{enumerate}
 \item The moment recurrence for $c_k$ and the closed moment--Thue--Morse formula
 for $F(a/2^n)$, including
 $F(1/4)=5/72$, $F(1/8)=1/288$,
 $F(3/8)=73/288$, and
 $F(5/16)=305857/2073600$.
 \item Direct sinc-product values of $a_k$ against the cosine reconstruction of
 $\Up$.
 \item The weighted moment residuals
 $\sum_{k<K}a_ko_k^{2m}$ for increasing $K$, together with rapid-decay tail
 bounds.
 \item The exact DCT formula \eqref{eq:A-F-DCT} against direct alias summation
 \eqref{eq:alias-identity}.
 \item The convergence table \cref{tab:A-convergence} and the rigorous bound
 \eqref{eq:alias-error-n}.
 \item The Galois conjugates at levels $n=2,3,4$, their rational power traces,
 and the orbit polynomials \eqref{eq:P2-poly}--\eqref{eq:P3-poly}.
 \item The single-alias factorization \eqref{eq:single-alias-factorization} for
 positive and negative alias indices.
 \item The observed sign law \eqref{eq:observed-A-sign} for every odd residue at
 levels $1\le n\le7$.
 \item The leading zero-jet approximation \eqref{eq:jet-error} for integers of
 several different dyadic valuations.
\end{enumerate}

\subsection{Exact arithmetic workflow}

The most stable exact path uses the even moments of $\Up$.  Set $c_0=1$ and
compute
\begin{equation}\label{eq:c-recurrence-verification}
 (2m+1)(2^{2m}-1)c_m
 =\sum_{j=0}^{m-1}\binom{2m+1}{2j}c_j.
\end{equation}
Then the repository's closed formula is
\begin{equation}\label{eq:closed-dyadic-verification}
 F\!\left(\frac a{2^n}\right)
 =\frac{2^{-\binom{n+1}{2}}}{n!}
 \sum_{h=0}^{a-1}\eps_h
 \sum_{k=0}^{\lfloor n/2\rfloor}
 \binom n{2k}(2a-2h-1)^{n-2k}c_k.
\end{equation}
Every quantity in \eqref{eq:closed-dyadic-verification} is rational.  It is
then inserted into \eqref{eq:A-F-DCT}; exact algebraic simplification is needed
only for the cosine values.

\subsection{A compact reproducibility listing}

\begin{lstlisting}[style=pseudo,language=Python,caption={Core exact computation used for the tables.}]
from fractions import Fraction
from math import comb, factorial


def tm_sign(k: int) -> int:
    return -1 if k.bit_count() % 2 else 1


def up_even_moments(max_k: int) -> list[Fraction]:
    c = [Fraction(1)]
    for m in range(1, max_k + 1):
        numerator = sum(
            Fraction(comb(2*m + 1, 2*j)) * c[j]
            for j in range(m)
        )
        denominator = (2*m + 1) * (2**(2*m) - 1)
        c.append(numerator / denominator)
    return c


def fabius_dyadic(a: int, n: int, c: list[Fraction]) -> Fraction:
    total = Fraction(0)
    for h in range(a):
        inner = sum(
            Fraction(comb(n, 2*k) * (2*a - 2*h - 1)**(n - 2*k))
            * c[k]
            for k in range(n//2 + 1)
        )
        total += tm_sign(h) * inner
    return total / Fraction(2**(n*(n + 1)//2) * factorial(n))

# A_{n,r} = 1/N - (2/N) sum_j F(j/N) cos(pi*r*j/N).
# Use exact Fraction values for F and a high-precision cyclotomic backend
# (or symbolic radicals for small n) for the cosine evaluation.
\end{lstlisting}

\begin{remark}[Reproducibility versus proof]
The computations are checks, not substitutes for the proofs.  Every central
identity in the report is established symbolically before numerical evidence is
invoked.  The code is included to make constants, orbit polynomials, and error
tables independently reproducible.
\end{remark}

\section{Lean formalization roadmap}\label{sec:lean}

The new theorem package separates naturally into finite algebra, elementary
analysis, and imported high-level asymptotics.  A low-risk formalization order is
as follows.

\subsection{Layer 1: finite Fourier algebra}

Formalize first the length-$2N$ DFT identities on a finite cyclic group:
\begin{enumerate}
 \item orthogonality of characters modulo $2N$;
 \item vanishing of all nonzero even folded modes;
 \item the inverse cosine transform \eqref{eq:inverse-DCT};
 \item finite Parseval and \eqref{eq:discrete-energy};
 \item the power-of-two Ramanujan sum \eqref{eq:Ramanujan-power-two};
 \item the rational trace formula \eqref{eq:Ramanujan-traces}.
\end{enumerate}
This layer uses only finite sums, rational dyadic values already available in
the project, roots of unity, and elementary complex algebra.

\subsection{Layer 2: Fourier-series extraction}

Use the established rapid decay of $\widehat\Up$ to formalize:
\begin{enumerate}
 \item the periodization $P$ and its complex Fourier coefficients;
 \item the reduction to odd modes and the cosine series
 \eqref{eq:cosine-spectrum};
 \item termwise differentiation to every order;
 \item the weighted Prouhet identities and polynomial quadrature;
 \item the derivative filtration \eqref{eq:filtration-above}--\eqref{eq:filtration-below};
 \item the exact spectral remainder \eqref{eq:F-rho-spectral}.
\end{enumerate}
The main analytic input is a reusable theorem that a Schwartz coefficient
sequence permits uniform termwise differentiation of a Fourier series.

\subsection{Layer 3: alias limits and estimates}

Formalize the infinite alias identity by applying the finite character filter to
an absolutely convergent Fourier series.  Then prove the selected-factor bound
and sum the shifted zeta tail.  The exact target statements are
\eqref{eq:alias-identity}, \eqref{eq:alias-error-general}, and
\eqref{eq:alias-error-n}.  No contour integration is needed.

\subsection{Layer 4: cyclotomic structure}

Two implementation choices are possible.
\begin{enumerate}
 \item Work concretely in $\C$ with a primitive root $\zeta_{2N}$ and prove
 invariance of symmetric polynomials under the explicit automorphisms.
 \item Work in an abstract cyclotomic extension, use the maximal real subfield
 only at the interface, and obtain rationality from fixed-field lemmas.
\end{enumerate}
The first route is likely shorter for the DCT formula; the second gives cleaner
Galois traces.  The orbit polynomial should be introduced without asserting it
is minimal.

\subsection{Layer 5: imported endpoint asymptotics}

Once \eqref{eq:F-rho-spectral} is formalized, the Lambert bridge is a one-line
rewrite using the project's existing endpoint theorem.  This is much easier than
formalizing a new saddle analysis.  A later project can address the genuinely
new localization problem posed in \cref{conj:Lambert-localization}.

\subsection{Indicative theorem interfaces}

\begin{lstlisting}[style=pseudo,caption={Schematic Lean theorem signatures.}]
-- Half-integer weighted Prouhet cancellation.
theorem halfIntegerSpectrum_evenMoment_zero (m : Nat) (hm : 0 < m) :
  HasSum (fun k => phi ((2*k+1 : Real) / 2) * (2*k+1)^(2*m)) 0

-- Exact folded DFT / alias identity.
theorem dyadicDCT_eq_aliasSum
    (n r : Nat) (hn : 0 < n) (hr : Odd r) (hrN : r < 2^n) :
  dyadicDCT n r = tsum (fun q : Int => phi (q * 2^n + r / 2))

-- Uniform certified recovery.
theorem dyadicDCT_sub_halfInteger_le
    (n r : Nat) (hn : 2 <= n) (hr : Odd r) (hrN : r < 2^n) :
  abs (dyadicDCT n r - phi (r / 2)) <= aliasError n

-- Rational Galois power trace.
theorem foldedPowerTrace_mem_rat
    (n m : Nat) : IsRat (sumOddResidues n (fun r => (dyadicDCT n r)^m))
\end{lstlisting}

\section{Further frontier problems}\label{sec:open}

\subsection{Exact signs of the folded coefficients}

\cref{conj:shifted-positivity} is the most immediate problem.  A proof would
turn the finite dyadic DCT into a second exact occurrence of the Thue--Morse sign
word, not merely an approximation to it.  Useful routes include:
\begin{enumerate}
 \item a quasi-periodic Green-kernel representation of
 $\sum_q\PhiU(q+\alpha)/(q+\alpha)^n$;
 \item variation-diminishing or total-positivity properties of dyadic samples of
 $\Up$;
 \item a recursion in $n$ for the quotient $A_{n,r}/\Gamma_{n,r}$;
 \item an integral representation with an explicitly positive kernel.
\end{enumerate}

\subsection{Minimal polynomials and stabilizers}

\begin{problem}\label{prob:minimal-polynomials}
Determine when $\mathcal P_n$ is the minimal polynomial of $A_{n,1}$ and when it
is a proper power.  More generally, classify all coincidences
$A_{n,r}=A_{n,s}$ and determine the discriminant and local factorization of the
orbit polynomial.
\end{problem}

The examples $n=2,3$ are irreducible over $\Q$.  The exact Galois action makes
this a finite arithmetic problem at each level, but no general primitivity
criterion is presently known.

\subsection{Arithmetic nature of the limiting amplitudes}

\begin{problem}\label{prob:amplitude-arithmetic}
Prove irrationality, transcendence, or nontrivial algebraic independence results
for
\[
 \PhiU(1/2),\ \PhiU(3/2),\ \PhiU(5/2),\ldots.
\]
The cyclotomic approximants \eqref{eq:q-approximant} come with unusually strong
error bounds, but their algebraic degrees and heights also grow.  Determine
whether the approximation quality can beat the relevant Liouville-type height
thresholds after sharp height estimates are included.
\end{problem}

\subsection{Spectral moments and nonlinear functionals of \texorpdfstring{$\Up$}{up}}

The trace limits \eqref{eq:trace-limit} define constants
$\mathfrak s_m=\sum_ka_k^m$.  The cases $m\ge3$ do not have an immediate
ordinary Parseval interpretation.

\begin{problem}\label{prob:nonlinear-traces}
Find direct integral, probabilistic, or operator-theoretic representations of
$\mathfrak s_m$ for $m\ge3$.  Determine whether the rational trace sequences
$p_m(n)$ satisfy recurrences or congruences inherited from the dyadic sample
array.
\end{problem}

\subsection{A genuine transform between the two Gaussian-binomial webs}

Formula \eqref{eq:q-approximant} composes the exact dyadic-value web with a
cyclotomic DCT.  The spectral determinant report constructs a different family
of Gaussian-binomial coefficients from finite q-Euler determinants.

\begin{problem}\label{prob:two-q-webs}
Construct an intrinsic matrix or incidence-algebra transform that takes the
Gaussian-binomial coefficients governing exact Fabius dyadic values directly
to the Gaussian-binomial coefficients governing the spectral q-Euler
determinants.  The DCT developed here supplies an explicit intermediate layer,
but not yet a closed transform between the two coefficient arrays.
\end{problem}

\subsection{Lambert spectral localization}

\begin{conjecture}[Lower-Lambert frequency localization]\label{conj:Lambert-localization}
Let $x\downarrow0$ and put $r(x)=\lambda(x)/x$ as in
\eqref{eq:lambda-W-def}.  There should exist windows
$\mathcal W_x$ centered on odd frequencies of order $r(x)$, with relative width
tending to zero on the logarithmic scale, such that the exact signed sum
\eqref{eq:F-rho-spectral} is asymptotically unchanged when restricted to
$\mathcal W_x$, after an appropriate multiscale grouping that preserves
Thue--Morse cancellation.

\smallskip\noindent\openstatus
\end{conjecture}

A naive termwise truncation is unlikely to work: the all-order moment identities
show that cancellations across scales are essential.  A successful formulation
may require dyadic blocks or a Mellin transform of the signed spectral measure
\eqref{eq:nu-measure}.

\subsection{A two-variable transform}

The corrected endpoint asymptotic is controlled by poles on the vertical
lattice $2\pi\ii\Z/\log2$.  The dyadic spectral tower has the same scale lattice,
and the folded DCT coefficients are finite projections of its half-integer
spectrum.

\begin{problem}\label{prob:two-variable-transform}
Construct a meromorphic two-variable transform whose residues simultaneously
produce
\begin{enumerate}
 \item the Gamma--zeta periodic endpoint function $\Psi$;
 \item the rational cyclotomic traces $p_m(n)$;
 \item the weighted Prouhet moment annihilation;
 \item the log-periodic heat trace of the dyadic spectral determinant.
\end{enumerate}
Such a transform would unify four presently separate manifestations of dyadic
self-similarity.
\end{problem}

\section{Conclusion}\label{sec:conclusion}

The half-integer values of the Rvachev sinc product sit at a previously
underdeveloped junction of the repository's main themes.  They are at once the
nonzero Fourier coefficients of the two-periodic up-function, the leading
constants in every integer-zero jet, and a Thue--Morse-signed sequence.  The
present report shows that these roles are not merely parallel observations.

Flatness at the center forces the sequence to solve an infinite positive-weight
Prouhet problem of every order.  Dyadic sampling folds the sequence into exact
finite discrete cosine transforms whose entries are rational Fabius values.
The finite $q=1/2$ Gaussian-binomial formula then turns each folded coefficient
into a completely explicit q-binomial--Thue--Morse--cyclotomic number.  An
elementary selected-factor estimate unfolds those coefficients with a uniform
superexponential error.  Galois conjugacy supplies rational orbit polynomials
and power traces, and those traces converge to nonlinear moments of the original
half-integer spectrum.  Splitting the first dyadic sinc factors reveals a finite
root-of-unity Thue--Morse product as the sign carrier of every alias class.
Finally, the all-order Prouhet cancellation rewrites the Fabius endpoint as an
exact spectral remainder beyond every algebraic order; the repository's
corrected Lambert--$W$ theorem then evaluates that same signed remainder
asymptotically.

The strongest immediate next theorem would be positivity of the residual
shifted spectral sum \eqref{eq:shifted-positivity}.  Beyond it lie arithmetic
questions about the cyclotomic tower, transcendence of the limiting amplitudes,
and a true lower-Lambert localization theorem.  The finite identities and error
bounds developed here are sufficiently elementary to support a staged Lean
formalization before those deeper questions are settled.

\appendix

\section{Corpus inventory}\label{app:inventory}

The same-day recursive audit snapshot counted fifty-seven TeX files: twenty-five in the living or imported documentation stratum and thirty-two frozen historical versions.  The current canonical frontier source is maintained by consolidation, so capitalization and the continued presence of processed draft paths can change after the snapshot. The list below records the path
relative to \path{Analysis/FabiusFunction/docs}. Historical items are grouped by
the archive directory; filenames follow the base name shown.

\subsection{Living documentation: 25 TeX files}

\begin{longtable}{@{}r p{0.88\textwidth}@{}}
\toprule
No. & Relative path \\
\midrule
\endfirsthead
\toprule
No. & Relative path \\
\midrule
\endhead
1 & \path{FabiusFunction_Mathematical_Glossary/FabiusFunction_Mathematical_Glossary.tex} \\
2 & \path{Fabius_Function_and_Rvachev_Up/Fabius_Function_and_Rvachev_Up.tex} \\
3 & \path{Non_Elementarity_of_the_Fabius_Function/Non_Elementarity_of_the_Fabius_Function.tex} \\
4 & \path{fabius_lean_walkthrough/fabius_lean_walkthrough.tex} \\
5 & \path{non-formalized-research-frontiers/non-formalized-research-frontiers.tex} \\
6 & \path{non-formalized-research-frontiers/drafts/Fabius_Thue_Morse_Formula_Atlas/Fabius_Thue_Morse_Formula_Atlas.tex} \\
7 & \path{non-formalized-research-frontiers/drafts/rvachev_up_fourier_decay/Rvachev_Up_Fourier_Decay-2/Rvachev_Up_Fourier_Decay.tex} \\
8 & \path{non-formalized-research-frontiers/drafts/rvachev_up_fourier_decay/Rvachev_Up_Fourier_Decay-3/Fourier_decay_of_Rvachev_up.tex} \\
9 & \path{non-formalized-research-frontiers/drafts/rvachev_up_fourier_decay/Rvachev_Up_Fourier_Decay_Comparative_Audit/Rvachev_Up_Fourier_Decay_Comparative_Audit.tex} \\
10 & \path{non-formalized-research-frontiers/drafts/rvachev_up_fourier_decay/Rvachev_Up_Fourier_Decay_Gentle_Guide/Rvachev_Up_Fourier_Decay_Gentle_Guide.tex} \\
11 & \path{non-formalized-research-frontiers/drafts/rvachev_up_fourier_decay/Rvachev_Up_Fourier_Decay_Second_Wave_Audit/Rvachev_Up_Fourier_Decay_Second_Wave_Audit.tex} \\
12 & \path{non-formalized-research-frontiers/drafts/rvachev_up_fourier_decay/asymptotic-decay-rate-of-an-infinite-product-of-sinc-functions/asymptotic-decay-rate-of-an-infinite-product-of-sinc-functions.tex} \\
13 & \path{non-formalized-research-frontiers/drafts/rvachev_up_fourier_decay/rvachev_up_fourier_decay-1/rvachev_up_fourier_decay.tex} \\
14 & \path{non-formalized-research-frontiers/drafts/rvachev_up_fourier_decay/rvachev_up_fourier_decay-4/rvachev_up_fourier_decay.tex} \\
15 & \path{non-formalized-research-frontiers/drafts/rvachev_up_fourier_decay/rvachev_up_fourier_decay-5/Rvachev_Up_Fourier_Decay_Final.tex} \\
16 & \path{non-formalized-research-frontiers/drafts/rvachev_up_fourier_decay/rvachev_up_fourier_decay-6/Rvachev_Up_Fourier_Decay_Final_Synthesis.tex} \\
17 & \path{non-formalized-research-frontiers/drafts/rvachev_up_fourier_decay/rvachev_up_fourier_decay-7/rvachev_up_fourier_decay_final.tex} \\
18 & \path{non-formalized-research-frontiers/drafts/rvachev_up_fourier_decay/rvachev_up_fourier_decay-8/rvachev_up_fourier_decay_final.tex} \\
19 & \path{papers/AIF_2017__67_2_637_0/AIF_2017__67_2_637_0.tex} \\
20 & \path{papers/AIF_2017__67_2_637_0/AIF_2017__67_2_637_0-corrected.tex} \\
21 & \path{papers/arXiv-1502.06738v2/Lacunary_products.tex} \\
22 & \path{papers/arXiv-1702.05442v1/09-Function.tex} \\
23 & \path{papers/arXiv-1702.06487v3/157-Arithmetic-v3.tex} \\
24 & \path{papers/arXiv-2211.13570v1/document.tex} \\
25 & \path{papers/ubiq15/ubiq15-corrected.tex} \\
\bottomrule
\end{longtable}

\subsection{Frozen archive: 32 TeX files}

\begin{longtable}{@{}r p{0.88\textwidth}@{}}
\toprule
No. & Relative path under \path{archive/} \\
\midrule
\endfirsthead
\toprule
No. & Relative path under \path{archive/} \\
\midrule
\endhead
1 & \path{early-notes/Fabius Asymptotic/Fabius Asymptotic.tex} \\
2 & \path{early-notes/K-fold summation over the signed Thue-Morse sequence/K-fold summation over the signed Thue-Morse sequence.tex} \\
3 & \path{standalone-studies/Non_Elementarity_of_the_Fabius_Function/Non_Elementarity_of_the_Fabius_Function.tex} \\
4 & \path{standalone-studies/Small_Argument_Asymptotics/Small_Argument_Asymptotics.tex} \\
5 & \path{standalone-studies/fabius-inverse-dyadic-closed-form/fabius-inverse-dyadic-closed-form.tex} \\
6 & \path{primary-exposition/Fabius_Function_and_Rvachev_Up/Fabius_Function_and_Rvachev_Up.tex} \\
7 & \path{primary-exposition/Fabius_Function_and_Rvachev_Up-2/Fabius_Function_and_Rvachev_Up.tex} \\
8 & \path{primary-exposition/Fabius_Function_and_Rvachev_Up-3/Fabius_Function_and_Rvachev_Up.tex} \\
9 & \path{primary-exposition/Fabius_and_Rvachev_Up-2/Fabius_and_Rvachev_Up.tex} \\
10 & \path{primary-exposition/Fabius_and_Rvachev_up/Fabius_and_Rvachev_up.tex} \\
11 & \path{primary-exposition/Fabius_function_and_up_function/Fabius_function_and_up_function.tex} \\
12 & \path{primary-supplements/Fabius_Function_Missing_Parts/Fabius_Function_Missing_Parts.tex} \\
13 & \path{primary-supplements/Fabius_Function_Missing_Parts-2/Fabius_Function_Missing_Parts.tex} \\
14 & \path{q-calculus/Fabius_q_Special_Functions/Fabius_q_Special_Functions.tex} \\
15 & \path{q-calculus/fabius-q-special-functions/fabius-q-special-functions.tex} \\
16 & \path{lean-walkthrough/Fabius_Function_Lean_Walkthrough/Fabius_Function_Lean_Walkthrough.tex} \\
17 & \path{lean-walkthrough/fabius_lean_walkthrough/fabius_lean_walkthrough.tex} \\
18 & \path{glossary/FabiusFunction_Mathematical_Glossary/FabiusFunction_Mathematical_Glossary.tex} \\
19 & \path{glossary/FabiusFunction_Mathematical_Glossary-2/FabiusFunction_Mathematical_Glossary.tex} \\
20 & \path{research-frontiers/Fabius_Dyadic_Asymptotic_Bridge/Fabius_Dyadic_Asymptotic_Bridge.tex} \\
21 & \path{research-frontiers/Fabius_Dyadic_Formulae_and_Alternative_Representations/Fabius_Dyadic_Formulae_and_Alternative_Representations.tex} \\
22 & \path{research-frontiers/Fabius_Dyadic_Formulae_to_Asymptotics/Fabius_Dyadic_Formulae_to_Asymptotics.tex} \\
23 & \path{research-frontiers/Fabius_Dyadic_q_Connections/Fabius_Dyadic_q_Connections.tex} \\
24 & \path{research-frontiers/Fabius_Thue_Morse_Convergence_Rate/Fabius_Thue_Morse_Convergence_Rate.tex} \\
25 & \path{research-frontiers/Fabius_Thue_Morse_Convergence_Rates/Fabius_Thue_Morse_Convergence_Rates.tex} \\
26 & \path{research-frontiers/Primary_Exposition_Gap_Register/Primary_Exposition_Gap_Register.tex} \\
27 & \path{research-frontiers/Repeated_Integration_and_Rvachev_Up/Repeated_Integration_and_Rvachev_Up.tex} \\
28 & \path{research-frontiers/Rvachev_Up_from_Repeated_Integration/Rvachev_Up_from_Repeated_Integration.tex} \\
29 & \path{research-frontiers/non-formalized-research-frontiers/non-formalized-research-frontiers.tex} \\
30 & \path{research-frontiers/Thue_Morse_Formula_Atlas-1/Thue_Morse_Formula_Atlas.tex} \\
31 & \path{research-frontiers/Thue_Morse_Formula_Atlas-2/Consolidated_Formula_Atlas_for_the_Thue_Morse_Sequence.tex} \\
32 & \path{research-frontiers/Thue_Morse_Formula_Atlas-3/Thue_Morse_Formula_Atlas.tex} \\
\bottomrule
\end{longtable}



\section{Selected source map}\label{app:sources}

\begin{thebibliography}{99}

\bibitem{proveit}
V.~Reshetnikov,
\emph{ProveIt: Analysis/FabiusFunction},
Lean~4 development and documentation corpus, living GitHub repository, 2026.
\url{https://github.com/VladimirReshetnikov/ProveIt/tree/main/Analysis/FabiusFunction}.

\bibitem{masterdoc}
V.~Reshetnikov,
\emph{The Fabius Function and Rvachev's Up-Function},
current main exposition in the \texttt{ProveIt} documentation corpus, 2026.
\url{https://github.com/VladimirReshetnikov/ProveIt/blob/main/Analysis/FabiusFunction/docs/Fabius_Function_and_Rvachev_Up/Fabius_Function_and_Rvachev_Up.tex}.

\bibitem{frontiervolume}
V.~Reshetnikov,
\emph{Non-formalized research frontiers for the Fabius--Rvachev development},
consolidated manuscript in the \texttt{ProveIt} repository, 2026.
\url{https://github.com/VladimirReshetnikov/ProveIt/tree/main/Analysis/FabiusFunction/docs/non-formalized-research-frontiers}.

\bibitem{spectralreport}
\emph{Dyadic Spectral Determinants for the Fabius--Rvachev System},
frontier report based on the complete recursive \texttt{ProveIt} corpus,
27 August 2026.

\bibitem{arithmeticrays}
\emph{Arithmetic Dyadic Rays of the Rvachev Fourier Product},
frontier report based on the complete recursive \texttt{ProveIt} corpus,
27 August 2026.

\bibitem{ariascompact}
J.~Arias de Reyna,
\emph{An infinitely differentiable function with compact support: Definition and properties},
arXiv:1702.05442, 2017.

\bibitem{ariasarith}
J.~Arias de Reyna,
\emph{Arithmetic of the Fabius function},
\emph{Integers} \textbf{18} (2018), Paper A51;
arXiv:1702.06487.

\bibitem{haugland}
J.~K.~Haugland,
\emph{Evaluating the Fabius function},
arXiv:1609.07999, 2016.

\bibitem{alloucheshallit1999}
J.-P.~Allouche and J.~Shallit,
\emph{The ubiquitous Prouhet--Thue--Morse sequence},
in \emph{Sequences and Their Applications}, Springer, 1999, 1--16.

\bibitem{alloucheshallitbook}
J.-P.~Allouche and J.~Shallit,
\emph{Automatic Sequences: Theory, Applications, Generalizations},
Cambridge University Press, 2003.

\bibitem{cloitre2026}
B.~Cloitre,
\emph{The Thue--Morse Transform},
arXiv:2604.06243v2, 2026.

\bibitem{sobolewski2019}
B.~Sobolewski,
\emph{Cyclotomic properties of polynomials associated with automatic sequences},
arXiv:1909.12052, 2019.

\bibitem{aistleitner2017}
C.~Aistleitner, R.~Hofer, and G.~Larcher,
\emph{On evil Kronecker sequences and lacunary trigonometric products},
\emph{Annales de l'Institut Fourier} \textbf{67} (2017), no.~2, 637--687.

\bibitem{toth2022}
L.~T\'oth,
\emph{Linear combinations of Dirichlet series associated with the Thue--Morse sequence},
arXiv:2211.13570, 2022.

\bibitem{corless}
R.~M.~Corless, G.~H.~Gonnet, D.~E.~G.~Hare, D.~J.~Jeffrey, and D.~E.~Knuth,
\emph{On the Lambert $W$ function},
\emph{Advances in Computational Mathematics} \textbf{5} (1996), 329--359.

\bibitem{washington}
L.~C.~Washington,
\emph{Introduction to Cyclotomic Fields}, 2nd ed.,
Graduate Texts in Mathematics 83, Springer, 1997.

\end{thebibliography}

\end{document}
